% s02-2-sl2.tex —— 2.2 节：G = SL_2 的情形（底本印刷页 30–35）

\section{\texorpdfstring{$G=\mathrm{SL}_2$ 的情形}{G=SL2 的情形}}
\label{sec:sl2}

本节考察 $G=\mathrm{SL}_2$ 的情形。无论在非分歧还是 Iwahori 分歧版本中，紧生成证明的
大部分概念性想法都可以在这个情形中看到。

\subsection{\texorpdfstring{$D\text{-}\mathrm{mod}(\Bun_{\mathrm{SL}_2})$ 的紧生成}{D-mod(SL2) 的紧生成}}
\label{sec:sl2-unram}

本小节复述 [DG15, 第 6 节] 中的讨论，它描述了 $G=\mathrm{SL}_2$ 时
$D\text{-}\mathrm{mod}(\Bun_G)$ 紧生成的论证。

局部 QCA 代数叠 $\Bun_G$ 表示 $X$ 上的主 $\mathrm{SL}_2$-丛，即行列式平凡、秩为 $2$
的向量丛 $\mathcal{M}$。$\Bun_G$ 是拟紧开子叠 $\Bun^{(\le m)}_G$ 的递增并，
$\Bun^{(\le m)}_G$ 表示这样的 $\mathrm{SL}_2$-丛 $\mathcal{M}$：它的每个线子丛的次数
至多为 $m$。

要看出 $m\ge \max(g-1,0)$ 时诸 $\Bun^{(\le m)}_G$ 余截断，只需说明
\[
\Bun^{(n)}_G=\Bun^{(\le n)}_G\setminus\Bun^{(\le n-1)}_G
\]
对 $n>\max(g-1,0)$ 截断。

现考虑分类短正合列
\[
0\to \mathcal{L}\to \mathcal{M}\to \mathcal{L}^{-1}\to 0
\]
的叠 $\Bun_B$，以及分类其中 $\dgr(\mathcal{L})=n$ 者的子叠 $\Bun^n_B$。有纤维方
\[
\begin{tikzcd}[column sep=huge,row sep=large]
\Bun^n_T \arrow[r,"i^-"] \arrow[d,"i"'] & \Bun^{-n}_B \arrow[d,"p^-"]\\
\Bun^n_B \arrow[r,"p"'] & \Bun^{(\le n)}_G,
\end{tikzcd}
\]
其中叠 $\Bun_T$ 分类 $X$ 上的线丛 $\mathcal{L}$，子叠 $\Bun^n_T$ 分类次数为 $n$ 的线丛，
\srcpage{30}
而映射 $i$ 与 $i^-$ 分别把 $\mathcal{L}$ 送到
\[
0\to \mathcal{L}\to \mathcal{L}\oplus \mathcal{L}^{-1}\to \mathcal{L}^{-1}\to 0
\]
与
\[
0\to \mathcal{L}^{-1}\to \mathcal{L}\oplus \mathcal{L}^{-1}\to \mathcal{L}\to 0.
\]

映射 $p$ 是到子叠 $\Bun^{(n)}_G\subset \Bun^{(\le n)}_G$ 上的同构，理由本质上是：
一个 $\mathrm{SL}_2$-丛至多只有一个次数为 $n$ 的线子丛。事实上，若 $\mathcal{L}$、
$\mathcal{L}'$ 是两个这样的子丛，则得到映射
\[
\mathcal{L}'\to \mathcal{M}\to \mathcal{L}^{-1}\in \Hom(\mathcal{L}',\mathcal{L}^{-1})
=H^0(X,\mathcal{L}'^{-1}\otimes \mathcal{L}^{-1})=0,
\]
因为 $\dgr(\mathcal{L}'^{-1}\otimes \mathcal{L}^{-1})=-2n<0$（反向亦然），故
$\mathcal{L}=\mathcal{L}'$。

映射 $i$ 满，因为每个带次数为 $n$ 的线子丛的 $\mathrm{SL}_2$-丛都分裂为直和：
$\Ext^1(\mathcal{L}^{-1},\mathcal{L})=H^1(X,\mathcal{L}^{\otimes 2})=0$，
而 $\dgr(\mathcal{L}^{\otimes 2})=2n>2g-2$。

映射 $p^-$ 光滑，因为 $p^-$ 在点
\[
0\to \mathcal{L}^{-1}\to \mathcal{M}\to \mathcal{L}\to 0
\]
处的微分的余纤维是 $H^1(X,\mathcal{L}^{\otimes 2})=0$。

最后说明映射 $i^-$。存在 $\Bun^n_T$ 上的向量丛 $\mathcal{V}$，其在
$\mathcal{L}\in \Bun^n_T$ 上的纤维等于
$\Ext^1(\mathcal{L},\mathcal{L}^{-1})=H^1(X,\mathcal{L}^{\otimes -2})$。
现记 $\mathrm{Pic}^n$ 为分类「$X$ 上次数为 $n$ 的线丛连同在某个固定点 $x\in X$ 处的
平凡化」的概形。考虑映射
\[
s\colon \mathrm{Pic}^n\to \Bun^n_T,
\]
\srcpage{31}
它忘掉平凡化的选取；它把 $\mathrm{Pic}^n$ 实现为 $\Bun^n_T$ 上的 $\mathbb{G}_m$-主丛。
沿 $s$ 拉回向量丛 $\mathcal{V}$，得到 $\mathrm{Pic}^n$ 上的向量丛。取零截面
\[
\mathrm{Pic}^n\to s^*\mathcal{V},
\]
再对 $\mathbb{G}_m$ 取商：
\[
\mathrm{Pic}^n/\mathbb{G}_m\to s^*\mathcal{V}/\mathbb{G}_m
\]
与映射
\[
i^-\colon \Bun^n_T\to \Bun^{-n}_B
\]
等同，故 $i^-$ 收缩。

\subsection{\texorpdfstring{$D\text{-}\mathrm{mod}(\Bun^{I}_{\mathrm{SL}_2})$ 的紧生成}{D-mod(ISL2) 的紧生成}}
\label{sec:sl2-iwahori}

现转向 Iwahori 分歧情形下的 $G=\mathrm{SL}_2$。

局部 QCA 代数叠 $\Bun^I_G$ 表示二元组 $(\mathcal{M},L)$，其中 $\mathcal{M}$ 是 $X$ 上的
主 $\mathrm{SL}_2$-丛，$L$ 是二维向量空间 $x^*\mathcal{M}$ 的一维子空间。当然有忘掉 $L$
的映射 $\Bun^I_G\to \Bun_G$。

叠 $\Bun^I_G$ 是拟紧开子叠 $\Bun^{I,(\le m)}_G$ 的递增并，后者是
$\Bun^{(\le m)}_G$ 沿 $\Bun^I_G\to \Bun_G$ 的基变换。

要看出 $m\ge\max(g-1,0)$ 时诸 $\Bun^{I,(\le m)}_G$ 余截断，只需说明
\[
\Bun^{I,(n)}_G=\Bun^{I,(\le n)}_G\setminus \Bun^{I,(\le n-1)}_G
\]
对 $n>\max(g-1,0)$ 截断。

叠 $\Bun^I_B$ 分类二元组
\[
\bigl(0\to \mathcal{L}\to \mathcal{M}\to \mathcal{L}^{-1}\to 0,\ L\bigr),
\]
而当 $n>\max(g-1,0)$ 时，
\[
\Bun^{I,n}_B\to \Bun^{I,(n)}_G
\]
\srcpage{32}
是同构。

叠 $\Bun^I_B$ 有两个自然的层。对 $\Bun^I_B$ 的一个给定点，可以问 $L$ 是否等于纤维
$x^*\mathcal{L}\subset x^*\mathcal{M}$。闭层记作 $\overline{\Bun}^{\,I}_B$，可等同于
$\Bun_B$；开层记作 $\opn{\Bun}^{I}_B$。\yizhuw{原文对闭层用「$\backslash\Bun$」这一
反斜杠前缀记号，与开层的「$^{\circ}\Bun$」相对；译本改用上划线，以免与集合差混淆。}

先看开层。有纤维方
\[
\begin{tikzcd}[column sep=huge,row sep=large]
\Bun^n_T \arrow[r,"i^-"] \arrow[d,"i"'] & \Bun^{-n}_B \arrow[d,"p^-"]\\
\opn{\Bun}^{I,n}_B \arrow[r,"p"'] & \Bun^{I,(\le n)}_G,
\end{tikzcd}
\]
其中 $i$ 把 $\mathcal{L}$ 送到
\[
\bigl(0\to \mathcal{L}\to \mathcal{L}\oplus \mathcal{L}^{-1}\to \mathcal{L}^{-1}\to 0,\
x^*\mathcal{L}^{-1}\bigr)\in \opn{\Bun}^{I,n}_B .
\]
断言 $i$ 满。设
$\bigl(0\to \mathcal{L}\to \mathcal{M}\to \mathcal{L}^{-1}\to 0,\ L\bigr)$ 是
$\opn{\Bun}^{I,n}_B$ 的一个点。已知 $\mathcal{M}$ 分裂为
$\mathcal{L}\oplus \mathcal{L}^{-1}$，故可把 $L$ 看作
$x^*\mathcal{L}\oplus x^*\mathcal{L}^{-1}$ 中的元素。设
\[
v\in x^*\mathcal{L},\quad w\in x^*\mathcal{L}^{-1}
\]
为 $x^*\mathcal{M}$ 的一组基。由于 $L\neq x^*\mathcal{L}$，存在形如
\[
\phi(v)=v,\qquad \phi(w)=cv+w
\]
的向量空间自同构 $\phi\in \Aut(x^*\mathcal{M})$，它把 $L$ 送到 $x^*\mathcal{L}^{-1}$。

下面说明 $\phi$ 可提升为 $\mathcal{M}$ 的自同构。有
\[
\End(\mathcal{M})=\Hom(\mathcal{L},\mathcal{L})\oplus \Hom(\mathcal{L}^{-1},\mathcal{L})
\oplus \Hom(\mathcal{L},\mathcal{L}^{-1})\oplus \Hom(\mathcal{L}^{-1},\mathcal{L}^{-1}).
\]
\srcpage{33}
$\Hom(\mathcal{L},\mathcal{L}^{-1})=H^0(X,\mathcal{L}^{\otimes -2})=0$，因为
$\dgr(\mathcal{L}^{\otimes -2})<0$。另一方面，
$\Hom(\mathcal{L}^{-1},\mathcal{L})=H^0(X,\mathcal{L}^{\otimes 2})$，而由
$\dgr(\mathcal{L}^{\otimes 2})>2g-1$ 知存在在 $x\in X$ 处不为零的整体截面 $f$；
不妨把它规范化，使其在 $x$ 处的值为 $c$。于是由
\[
\Phi=(\mathrm{id}_{\mathcal{L}},f,0,\mathrm{id}_{\mathcal{L}}^{-1})
\in \Hom(\mathcal{L},\mathcal{L})\oplus \Hom(\mathcal{L}^{-1},\mathcal{L})
\oplus \Hom(\mathcal{L},\mathcal{L}^{-1})\oplus \Hom(\mathcal{L}^{-1},\mathcal{L}^{-1})
\]
给出的 $\mathcal{M}$ 的自同构限制为 $\phi$。

于是 $\Phi$ 给出同构
\[
\bigl(0\to \mathcal{L}\to \mathcal{M}\to \mathcal{L}^{-1}\to 0,\ L\bigr)
\xrightarrow{\ \sim\ }
\bigl(0\to \mathcal{L}\to \mathcal{M}\to \mathcal{L}^{-1}\to 0,\ x^*\mathcal{L}^{-1}\bigr),
\]
而后者显然落在 $i$ 的像中。

映射 $p^-$ 光滑，因为 $p^-$ 在点
\[
0\to \mathcal{L}^{-1}\to \mathcal{M}\to \mathcal{L}\to 0
\]
处的微分的余纤维是 $H^1(X,\mathcal{L}^{\otimes 2}(-x))=0$，因为
$\dgr(\mathcal{L}^{\otimes 2}(-x))\ge 2g-1$。

映射 $i^-$ 与非分歧情形中的映射相同，前面已说明它收缩。
\srcpage{34}

闭层也有纤维方：
\[
\begin{tikzcd}[column sep=huge,row sep=large]
\Bun^n_T \arrow[r,"i^-"] \arrow[d,"i"'] & \opn{\Bun}^{I,-n}_B \arrow[d,"p^-"]\\
\Bun^n_B \arrow[r,"p"'] & \Bun^{I,(\le n)}_G .
\end{tikzcd}
\]
这里 $i$ 与非分歧情形中的映射相同，故已知它满。

映射 $p^-$ 可写成复合
\[
\opn{\Bun}^{I,-n}_B\to \Bun^{I,-n}_B\to \Bun^{I,(\le n)}_G .
\]
第一个映射是开嵌入，第二个映射是光滑映射 $\Bun^{-n}_B\to \Bun^{(\le n)}_G$ 的基变换，
故 $p^-$ 光滑。

下面说明映射 $i^-$。存在 $\Bun^n_T$ 上的向量丛 $\mathcal{V}$，其在
$\mathcal{L}\in \Bun^n_T$ 上的纤维等于 $H^1(X,\mathcal{L}^{\otimes -2}(-x))$。
同样可沿 $s\colon \mathrm{Pic}^n\to \Bun^n_T$ 拉回 $\mathcal{V}$。这一次取零截面
\[
\mathrm{Pic}^n\to s^*\mathcal{V},
\]
并对 $\mathbb{G}_m$ 取商：
\[
\mathrm{Pic}^n/\mathbb{G}_m\to s^*\mathcal{V}/\mathbb{G}_m
\]
与
\[
i^-\colon \Bun^n_T\to \opn{\Bun}^{I,-n}_B
\]
等同，故后者收缩。
\srcpage{35}
